\chapter{Fundamentos teóricos}
\label{cap:fundamentos}

Este capítulo introduce los conceptos necesarios para comprender el trabajo: los
fundamentos del Análisis Topológico de Datos, con énfasis en la homología
persistente, y una breve descripción de los modelos de aprendizaje profundo
empleados como referencia.

\section{Análisis Topológico de Datos}
El TDA parte de la idea de que los datos poseen una \emph{forma} y que dicha
forma contiene información relevante. Formalmente, se construye a partir de los
datos una familia de espacios topológicos y se estudia cómo evolucionan sus
invariantes.

\subsection{Complejos simpliciales y filtraciones}
\begin{definicion}[Complejo simplicial]
Un \emph{complejo simplicial} $K$ es un conjunto de símplices (vértices, aristas,
triángulos, tetraedros, \dots) cerrado bajo la operación de tomar caras y tal
que la intersección de dos símplices es una cara de ambos.
\end{definicion}

Para analizar datos se construye una \emph{filtración}, es decir, una familia
anidada de complejos simpliciales
\[
  \emptyset = K_0 \subseteq K_1 \subseteq \cdots \subseteq K_n = K,
\]
indexada por un parámetro de escala. En imágenes, una filtración habitual es la
\emph{filtración por sublevel sets} de la intensidad: se van incluyendo los
píxeles cuyo valor es menor o igual que un umbral creciente.

\subsection{Homología y números de Betti}
La homología asocia a cada complejo una secuencia de grupos cuyos rangos, los
\emph{números de Betti} $\beta_k$, cuentan las características topológicas de
dimensión $k$:
\begin{itemize}
  \item $\beta_0$: número de \textbf{componentes conexas}.
  \item $\beta_1$: número de \textbf{agujeros} o ciclos unidimensionales.
  \item $\beta_2$: número de \textbf{cavidades} o huecos tridimensionales.
\end{itemize}

\subsection{Homología persistente}
La \emph{homología persistente} rastrea el nacimiento (\textit{birth}) y la
muerte (\textit{death}) de las características topológicas a lo largo de la
filtración. Una característica que persiste durante un amplio rango de escalas se
considera una propiedad estructural significativa, mientras que las de vida corta
suelen corresponder a ruido.

\begin{definicion}[Diagrama de persistencia]
El \emph{diagrama de persistencia} es un multiconjunto de puntos
$(b_i, d_i) \in \mathbb{R}^2$ con $b_i \le d_i$, donde cada punto representa una
característica topológica que nace en la escala $b_i$ y muere en $d_i$. La
\emph{persistencia} de la característica es $d_i - b_i$.
\end{definicion}

La misma información puede representarse como un \emph{código de barras}
(\textit{barcode}), donde cada característica es un intervalo $[b_i, d_i]$.

\subsection{Estabilidad y vectorización}
Una propiedad fundamental es la \textbf{estabilidad}: pequeñas perturbaciones en
los datos producen pequeñas variaciones en el diagrama de persistencia (medidas
con la distancia \textit{bottleneck} o de Wasserstein). Esto explica la robustez
del TDA frente al ruido.

Para poder utilizar los diagramas como entrada de algoritmos de aprendizaje
automático es necesario \emph{vectorizarlos}. Entre las representaciones más
utilizadas destacan:
\begin{itemize}
  \item \textbf{Imágenes de persistencia} (\textit{persistence images}):
        transforman el diagrama en una imagen (matriz) de tamaño fijo.
  \item \textbf{Persistence landscapes}: funciones a trozos que resumen el
        diagrama.
  \item \textbf{Curvas de Betti} y estadísticos de persistencia.
\end{itemize}

\section{Aprendizaje profundo (referencia)}
Como referencia comparativa se emplean redes neuronales convolucionales (CNN).

\subsection{Redes neuronales convolucionales}
Una CNN aprende, mediante capas convolucionales, jerarquías de características
espaciales a partir de las imágenes. Combinada con capas de \textit{pooling} y
capas totalmente conectadas, permite abordar tareas de clasificación de imagen.

\subsection{Transfer learning}
Dado el tamaño limitado de los conjuntos de datos médicos, se recurre al
\textit{transfer learning}: se parte de una red preentrenada en un gran corpus
(p.\,ej. ImageNet) y se ajusta (\textit{fine-tuning}) a la tarea de mamografía.
En este trabajo se utiliza como referencia una arquitectura de la familia
DenseNet/ResNet.

\begin{ejemplo}[Intuición del enfoque topológico en mamografía]
Una agrupación de microcalcificaciones puede caracterizarse por el número de
componentes conexas ($\beta_0$) y su patrón de aparición a lo largo de la
filtración de intensidad, mientras que la estructura de una masa espiculada
puede reflejarse en la aparición y desaparición de ciclos ($\beta_1$). El TDA
codifica estos patrones de forma compacta y robusta.
\end{ejemplo}
